\begin{table}[htbp]
\centering
\footnotesize
\setlength{\tabcolsep}{4pt}
\sloppy
\caption{Transition episodes by hazard, county panel 2011--2023}
\label{tab:e2t1}
\begin{tabularx}{\textwidth}{>{\raggedright\arraybackslash}X r r r r r r}
\toprule
Hazard & Onset & Persist & Exit & Calm & County-years & Counties \\
\midrule
Extreme drought & 511 & 175 & 705 & 36,253 & 37,644 & 603 \\
Extreme heat & 903 & 8,125 & 1,033 & 27,583 & 37,644 & 1,016 \\
Extreme cold & 1,032 & 5,485 & 1,181 & 29,946 & 37,644 & 834 \\
\bottomrule
\end{tabularx}
\begin{minipage}{\textwidth}\vspace{0.5em}\footnotesize
A county-year is classified by its position in an exposure episode: onset (first year of a spell), persistence (a continuing spell), exit (the first year after a spell ends), or calm. The counts show how much support each transition has, which is what limits the horizons the dynamic estimates can reach: drought spells are rare and short, while heat spells persist.
\end{minipage}
\end{table}
